We calculated 77 image-part/timing jobs for 63 events.
63 events have defined style ratios under the stated primary assumptions.
The calculated range is 0.03879--463.9, from G23
to B10. Signature selection used the newly calculated avulsion
areas. It selected brightness for 57 events and
wetness for 6 events.

Every deposited event mask is named for only one spectral signature. We therefore
apply both filters to the same supplied geometry, retaining the mask's source name
only as provenance. This tests numerical signature selection on fixed geometry;
it cannot recreate a separately redrawn mask for the alternative signature.

The embedded image calendar labels disagree with 33 selected metadata
labels across 6 jobs. Native band selections
remain unchanged. The file \texttt{temporal\_labels.csv} records every assignment.
The unused-input ledger includes B5, which lacks a corresponding selected-window
record, and extra B10 tiles. No excluded measured result is used to fill these gaps.

The overlap audit identifies 15 site/signature/role combinations with
duplicate detected pixels. The primary event counts take the spatial union on
aligned native grids, so a physical pixel is counted once. A separate summed-part
calculation quantifies the effect of simply adding tile counts. Mosaic colors use
the maximum per-part detection count in overlaps; heterogeneous windows prevent
interpreting that color as one uniform observation period.

Across the 146 primary image-part/role windows, the fraction of
mask pixels observed at least once ranges from 91.59 to 100 percent;
22 windows have some never-observed pixels. This cumulative
coverage does not imply full coverage in every year. Missing pixels are not
filled, and the ratio is not corrected by an assumed missing-area multiplier.

We independently annotated 4 prespecified illustrative events
(B11, B12, V11 and V15). These are single-date, 30 m observations, not a new
63-event morphology survey. The elevation exercise produced 2 positive
slopes meeting the explicit reach-length and sign checks. Widths are wet-pixel
cross-section estimates with a native pixel scale of 30 m. No mountain-front
distance was measured. Consequently, the full distance relationship and the
published population distributions remain unreproduced.
